\SourceLeaf{322}%
\par\noindent\hspace*{\parindent} Zuvörderst ergibt sich aus der Verbindung der Gleichungen (1), (2), (3), dass \(qq+\alpha\alpha ll-uu-ss\) eine Grösse vierter Ordnung ist, daher man anstatt (2) auch schreiben kann
\[
l=\frac{s\mathord{.}\sin U}{\alpha\mathord{.}\cos Q}(1-\tfrac1{24}qq+\tfrac1{24}uu)
\]
oder wenn man nach (8), (9) und (10)
\[
s=\frac{r\cos\theta^2}{a\cos\varphi},\qquad \alpha\cos Q=\frac{\cos\theta\cos P}{\cos\varphi}
\]
schreibt,
\[
l=\frac{r\cos\theta\sin U}{a\cos P}(1-\tfrac1{24}qq+\tfrac1{24}uu)
\]
Es wird ferner \(\frac{\cos\theta}{\cos P}=\frac{\BareRoot{(1-ee\sin P^2)}}{\cos P}\) vermittelst der Gleichung (4) und durch eine leichte Rechnung entwickelt in
\[
\frac{\cos\theta}{\cos P}=\frac{\BareRoot{(1-ee\sin B^2)}}{\cos B}(1+\frac{3ee\sin P^2}{8(1-ee\sin P^2)}\mathord{.}qq)
\]
was bis auf die vierte Ordnung ausschl. genau ist. Wir haben daher, wenn zugleich \(T\) für \(U\) geschrieben wird, gemäss der Gleichung (6),
\[
l=\frac{r\BareRoot{(1-ee\sin B^2)}\mathord{.}\sin T}{a\cos B}(1-\frac{1-10ee\sin P^2}{24(1-ee\sin P^2)}\mathord{.}qq+\tfrac1{24}uu)
\]
Nachdem in dem eingeklammerten Theile noch substituirt ist \(\frac{\cos\varphi}{\BareRoot{(1-ee\sin P^2)}}\mathord{.}b\) für \(q\), sodann \(t\) für \(u\) und endlich \(B\) für \(P\), was alles, nach Gleichung (5), (7), (4), wie man leicht sieht, geschehen kann, ohne den Grad der Genauigkeit zu vermindern, und wenn wir ausserdem, zur Abkürzung,
\[
\BareRoot{(1-ee\sin B^2)}=k
\]
schreiben, so erhalten wir (I)
\[
l=\frac{kr\sin T}{a\cos B}(1-\frac{(1-10ee\sin B^2)\cos\varphi^2}{24k^4}\mathord{.}bb+\tfrac1{24}tt)
\]

\Article{24}
Auf ähnliche Weise verwandelt sich Gleichung (1) in
\[
q=s\cos U(1-\tfrac1{12}qq+\tfrac1{12}ss+\tfrac18uu)
\]
und daher Gleichung (5) in
